% ============================================================
%  DebriSolver Paper — LaTeX Version
%  Saudi Space Agency | DebriSolver Competition | SDC 2026
%  King Abdulaziz University
% ============================================================

\documentclass[12pt, a4paper]{article}

% ── Packages ──────────────────────────────────────────────
\usepackage[T1]{fontenc}
\usepackage[utf8]{inputenc}
\usepackage{lmodern}
\usepackage{geometry}
\geometry{a4paper, margin=2.5cm}

\usepackage{amsmath}
\usepackage{amssymb}
\usepackage{bm}

\usepackage{graphicx}
\graphicspath{{images/}}

\usepackage{float}
\usepackage{caption}
\usepackage{subcaption}

\usepackage{booktabs}
\usepackage{array}
\usepackage{longtable}
\usepackage{multirow}

\usepackage{hyperref}
\hypersetup{
    colorlinks=true,
    linkcolor=blue,
    urlcolor=blue,
    citecolor=blue,
    pdftitle={Learning Conjunction Dynamics: A Self-Supervised Approach to Satellite Collision Risk Assessment},
    pdfauthor={Ahmad Alharbi, Abdulelah Mojelad, Hamzah Alharbi, Khalid Alsadoon, Mohamedhakim Hassan}
}

\usepackage{parskip}
\usepackage{setspace}
\setstretch{1.15}

\usepackage{enumitem}
\usepackage{xcolor}
\usepackage{titlesec}

% Section formatting
\titleformat{\section}{\large\bfseries}{\thesection.}{0.5em}{}
\titleformat{\subsection}{\normalsize\bfseries}{\thesubsection}{0.5em}{}
\titleformat{\subsubsection}{\normalsize\itshape\bfseries}{\thesubsubsection}{0.5em}{}

% ── Title ─────────────────────────────────────────────────
\title{%
    \LARGE\bfseries Learning Conjunction Dynamics:\\
    A Self-Supervised Approach to Satellite Collision\\
    Risk Assessment
}

\author{
    Ahmad Alharbi\textsuperscript{1} \quad
    Abdulelah Mojelad\textsuperscript{2} \quad
    Hamzah Alharbi\textsuperscript{3} \quad
    Khalid Alsadoon\textsuperscript{4} \quad
    Mohamedhakim Hassan\textsuperscript{5}
    \\[0.8em]
    \small\textsuperscript{1}Department of Aerospace Engineering, Faculty of Engineering,\\
    King Abdulaziz University, Jeddah 23835\\
    \href{mailto:ahmadharbi157@hotmail.com}{ahmadharbi157@hotmail.com}\\[0.4em]
    \small\textsuperscript{2}Department of Computer Science, Faculty of Computing and Information Technology,\\
    King Abdulaziz University, Jeddah 23826\\
    \href{mailto:abdulellah.mazen@gmail.com}{abdulellah.mazen@gmail.com}\\[0.4em]
    \small\textsuperscript{3}Department of Aerospace Engineering, Faculty of Engineering,\\
    King Abdulaziz University, Jeddah 23353\\
    \href{mailto:hamzahawad789@gmail.com}{hamzahawad789@gmail.com}\\[0.4em]
    \small\textsuperscript{4}Department of Aerospace Engineering, Faculty of Engineering,\\
    King Abdulaziz University, Jeddah 23235\\
    \href{mailto:eng.khalid.alsadoon@gmail.com}{eng.khalid.alsadoon@gmail.com}\\[0.4em]
    \small\textsuperscript{5}Department of Aerospace Engineering, Faculty of Engineering,\\
    King Abdulaziz University, Jeddah 23241\\
    \href{mailto:mohamedhakim1901@gmail.com}{mohamedhakim1901@gmail.com}
}

\date{}

% ============================================================
\begin{document}
\maketitle
\thispagestyle{empty}

% ── Abstract ──────────────────────────────────────────────
\begin{abstract}
The rapid increase in tracked space objects has introduced significant operational challenges for satellite operators, particularly in distinguishing credible collision threats from the large volume of low-quality conjunction alerts generated by space surveillance networks. This study introduces a self-supervised deep learning method for assessing the credibility of conjunction alerts, addressing the core limitation that actual satellite collisions are too infrequent to provide labeled training data for supervised classification. A bidirectional gated recurrent unit architecture with Monte Carlo Dropout was trained on 185,511 conjunction data messages collected from January to June 2024. The model learned to predict the evolution of collision probability, miss distance, and covariance parameters from historical conjunction data message sequences without relying on human-defined threat labels. On a test set of 6,411 unique conjunction events, the model achieved a mean absolute error of 0.403 for collision probability prediction and a log-scaled collision probability $R^2$ of 0.649, demonstrating effective learning across an eight-order-of-magnitude probability range from $10^{-10}$ to $10^{-2}$. A four-quadrant risk classification framework that integrates model-derived threat scores with Monte Carlo Dropout uncertainty estimates produces operationally realistic distributions. This framework identified 4 percent of events as requiring immediate attention and safely deprioritized 56.5 percent as routine conjunctions, resulting in a 96 percent reduction in urgent alerts. The approach supports real-time processing, with 25-minute training on standard CPU hardware and millisecond-scale inference, making it suitable for operational deployment in space traffic management workflows. These findings demonstrate that self-supervised learning with uncertainty quantification can effectively address conjunction screening challenges without requiring collision outcome labels or fixed probability thresholds.

\bigskip
\noindent\textbf{Keywords:} Self-Supervised Learning; Satellite Conjunction Screening; Monte Carlo Dropout; Uncertainty Quantification; Space Traffic Management.
\end{abstract}

\newpage

% ── Table of Contents ──────────────────────────────────────
\tableofcontents
\newpage

% ============================================================
\section{Introduction}

The sustainability of operations in Low Earth Orbit (LEO) faces significant threats due to increasing orbital congestion, primarily caused by the rapid deployment of mega-constellations and the heightened risk of catastrophic debris generation \cite{ref1}. This congestion imposes a large volume of Conjunction Data Messages (CDMs) on satellite operators, which requires timely and precise avoidance maneuvers. This operational context reveals a critical limitation in traditional risk assessment approaches: reliance on fixed thresholds or deterministic models often results in excessive false-positive alerts, leading to unnecessary propellant expenditure and compromising both mission sustainability and profitability \cite{ref2}.

\subsection{The Operational Challenge and Technical Gap}

The overwhelming volume of CDMs has rendered manual processing by human operators unsustainable, necessitating advanced automation. This challenge is intensified by the vast number of debris fragments in orbit and the persistent risk of the Kessler Syndrome, in which a single collision can trigger a cascade of additional debris \cite{ref2}. A central operational challenge is the lack of a robust mechanism to assess prediction reliability before executing costly maneuvers.

Traditional methods, which depend on fixed probabilistic thresholds or deterministic models, remain highly vulnerable to data instability and uncertainties in covariance measurements. This vulnerability compels operators to choose between accepting the risk of a catastrophic false negative, such as a missed collision, or incurring the guaranteed cost of a resource-depleting false positive, such as an unnecessary maneuver. The consequent depletion of the satellite's propellant reserves directly undermines mission sustainability.

To address this challenge, this study applies Artificial Intelligence and Machine Learning (AI/ML) techniques to develop a scalable, confidence-weighted operational solution.

\subsection{Rationale, Solution, and Technical Innovations}

To achieve the required reliability, this research introduces the \textbf{Self-Supervised CDM Trajectory Learning System}, a novel framework designed to convert raw CDM time-series data into confidence-weighted operational directives. The methodology is based on two core principles: addressing the lack of ground truth labels and quantifying prediction uncertainty.

\begin{itemize}
    \item \textbf{Self-Supervised Trajectory Learning:} To address the challenge posed by the extreme rarity of actual collision outcomes, which makes standard supervised learning impractical, a self-supervised learning paradigm is employed. The model is trained to predict the parameters of the subsequent CDM ($\text{CDM}_n$) in a sequence, using its historical data ($\text{CDM}_1$ to $\text{CDM}_{n-1}$). This prediction task leverages a robust Bidirectional Gated Recurrent Unit (BiGRU) architecture, a variant of Recurrent Neural Networks (RNNs), to capture the generalized trajectory dynamics of conjunction evolution. This approach enables the model to learn valid dynamics directly from data patterns.

    \item \textbf{Built-in Uncertainty Quantification:} Monte Carlo Dropout is applied during inference to generate a continuous Model Confidence Level. This measure is essential for distinguishing between genuine threats and data instability.
\end{itemize}

The primary objective of this research is to validate the effectiveness of confidence-weighted risk assessment in optimizing operational resource efficiency. The main conclusion is that the uncertainty-aware framework delivers superior operational efficiency. The model effectively filters the majority of screened events into the \textbf{SAFELY IGNORE} category, providing operators with the assurance required to conserve propellant while maintaining high-confidence identification of the limited cases necessitating immediate intervention. By converting probabilistic risk into clear, confidence-supported operational directives, this solution directly contributes to the long-term sustainability of the LEO environment.

% ============================================================
\section{Methodology}

Our method integrates robust data handling, trajectory-based sequence learning, bidirectional recurrent neural network architecture, and Monte Carlo uncertainty quantification into a unified operational framework. Each methodological component is explained in detail, from reading raw CDMs files through the final four-quadrant risk classification scheme. This will make the workflow easy to repeat and understand the proposed solution.

\subsection{Overview and Approach}

A self-supervised deep learning architecture is presented for satellite conjunction risk analysis. The model learns conjunction dynamics directly from historical CDM sequences, quantifying both predictive uncertainty and assessment confidence without relying on the threshold-based rules or probabilistic components common to traditional approaches. The framework comprises five core stages:

\begin{enumerate}
    \item Robust parsing and enrichment of raw CDM data in Key-Value Notation (KVN) format.
    \item Feature engineering and self-supervised sequence construction.
    \item Bidirectional GRU training with Monte Carlo (MC) Dropout.
    \item Behavior-derived threat and confidence scoring.
    \item Four-quadrant risk classification.
\end{enumerate}

\subsection{Data Acquisition and Preprocessing}

The dataset contains 185,511 CDMs, covering global conjunction screenings from 1 January to 30 June 2024, sourced from Aldoria and the SSA DebrisSolver competition partnership. Due to severe timestamp corruption in the supplied CSV (approximately 60\% unreadable), a custom Python parser was developed to process the original KVN files directly. The parser extracts 11 features per CDM:

\begin{itemize}
    \item Collision probability $P_c$
    \item $\log_{10}(P_c)$
    \item Time to Closest Approach (TCA) in hours
    \item Combined Radial-Transverse-Normal (RTN) positional uncertainty $(c_{r,r} + c_{t,t} + c_{n,n})$
    \item Relative speed magnitude
    \item Three relative position components in the primary object's RTN frame
    \item Three relative velocity components in the primary object's RTN frame
\end{itemize}

The custom parsing approach achieved 99.95\% timestamp recovery, retaining 185,415 valid CDMs out of the original 185,511. Features were standardized using a \texttt{StandardScaler} fitted on the training set. Prior to normalization, miss distance was clipped to the range $[1,\, 10^5]$\,m and $P_c$ to $[10^{-10},\, 10^{-2}]$.

\subsection{Self-Supervised Sequence Learning}

Conjunction events were grouped by unique concatenated NORAD IDs and sorted chronologically. Events containing $\geq 2$ CDMs were formed into input--target pairs. Given a historical sequence up to timestep $n-1$, the task is defined as \cite{ref3}:

\begin{equation}
    \text{Input: } \text{CDM}_1, \ldots, \text{CDM}_{n-1} \;\rightarrow\; \text{Target: } \text{CDM}_n
    \quad \text{for all } n \in \{2, \ldots, N\}
\end{equation}

Using past CDMs, the model predicts the feature vector at timestep $n$ via the function:

\begin{equation}
    f:\; \mathbb{R}^{(n-1)\times 11} \;\rightarrow\; \mathbb{R}^{11}
\end{equation}

The maximum sequence length was fixed at 20 timesteps, which covers 86.4\% of events directly (90\% of events contain 13 or fewer CDMs). Events exceeding 20 CDMs generated multiple overlapping windows of 20, while shorter events were zero-padded to the fixed size and masked during training. An event-level random 80/10/10 split yielded:

\begin{itemize}
    \item \textbf{Training:} 77,989 sequences (51,287 events)
    \item \textbf{Validation:} 9,598 sequences (6,411 events)
    \item \textbf{Test:} 9,637 sequences (6,411 events)
\end{itemize}

\subsection{Neural Network Architecture and Training}

\begin{figure}[H]
    \centering
    \includegraphics[width=0.92\textwidth]{fig1_model_architecture.png}
    \caption{The model architecture features two stacked Bidirectional GRU layers followed by dense layers, incorporating Batch Normalization and Monte Carlo Dropout for uncertainty quantification.}
    \label{fig:architecture}
\end{figure}

The proposed solution utilizes a stacked bidirectional GRU network \cite{ref4} capable of processing variable-length sequences in both forward and backward temporal directions, as illustrated in Figure~\ref{fig:architecture}.

The architecture begins with a masking layer to handle padded inputs, followed by:

\begin{enumerate}
    \item A first \textbf{Bidirectional GRU} layer with 128 units in each direction (256 total). Outputs are concatenated, normalized via Batch Normalization, and regularized with Dropout ($p = 0.3$).
    \item A second \textbf{Bidirectional GRU} layer with 64 units per direction (128 total), also followed by concatenation, Batch Normalization, and Dropout ($p = 0.3$).
    \item Two fully connected \textbf{Dense} layers (64 and 32 units, ReLU activation), each followed by Dropout ($p = 0.3$).
    \item A final output projection to an \textbf{11-feature vector}.
\end{enumerate}

The model with \textbf{244,171 parameters} was trained to minimize the Weighted Mean Squared Error (wMSE), assigning higher weights to critical features: $\log_{10}(P_c)$ (2.0), miss distance (1.5), and uncertainty features (1.2). Training used the Adam optimizer with $\text{lr} = 0.001$, batch size 64, maximum 150 epochs, with early stopping (patience = 20) and learning rate reduction on plateau.

Model performance on the held-out test set (10\%) is quantified using three complementary regression metrics \cite{ref5, ref6}. These metrics are defined by the following equations:

\subsubsection{Weighted Mean Squared Error (wMSE)}

Batch loss function:

\begin{equation}
    L = \frac{1}{B} \sum_{i=1}^{B} \sum_{j=1}^{11} w_j \left(\hat{y}_{ij} - y_{ij}\right)^2
    \label{eq:batch_loss}
\end{equation}

Weighted Mean Squared Error:

\begin{equation}
    \text{wMSE} = \frac{1}{N} \sum_{i=1}^{N} \sum_{j=1}^{11} w_j \left(\hat{y}_{ij} - y_{ij}\right)^2
    \label{eq:wmse}
\end{equation}

\subsubsection{Mean Absolute Error (MAE)}

Per-feature and overall MAE, calculated as:

\begin{equation}
    \text{MAE}_j = \frac{1}{N} \sum_{i=1}^{N} \left| y_{ij} - \hat{y}_{ij} \right|, \qquad
    \text{MAE}_{\text{overall}} = \frac{1}{N \times 11} \sum_{i=1}^{N} \sum_{j=1}^{11} \left| y_{ij} - \hat{y}_{ij} \right|
    \label{eq:mae}
\end{equation}

MAE for collision probability:

\begin{equation}
    \text{MAE}_{P_c} = \frac{1}{N} \sum_{i=1}^{N} \left| P_{c_i}^{\text{pred}} - P_{c_i}^{\text{true}} \right|
    \label{eq:mae_pc}
\end{equation}

\subsubsection{Coefficient of Determination ($R^2$)}

Per-feature $R^2$:

\begin{equation}
    R_j^2 = 1 - \frac{\displaystyle\sum_{i=1}^{N} \left(y_{ij} - \hat{y}_{ij}\right)^2}
                     {\displaystyle\sum_{i=1}^{N} \left(y_{ij} - \bar{y}_j\right)^2}
    \label{eq:r2_j}
\end{equation}

Overall $R^2$:

\begin{equation}
    R^2 = \frac{1}{11} \sum_{j=1}^{11} R_j^2
    \label{eq:r2_overall}
\end{equation}

\subsubsection{Uncertainty Quantification}

Uncertainty is quantified via Monte Carlo (MC) Dropout at inference — 50 stochastic forward passes ($T = 50$, $p = 0.3$) \cite{ref7}. The final prediction is the mean across passes and the epistemic (model) uncertainty is the standard deviation. Equations~\eqref{eq:mc_mean} and~\eqref{eq:mc_std} calculate mean and standard deviation \cite{ref8, ref9, ref10}:

\begin{equation}
    \hat{y} = \frac{1}{T} \sum_{i=1}^{T} f\!\left(\mathbf{x};\, \bm{w},\, \mathbf{d}_i\right)
    \label{eq:mc_mean}
\end{equation}

\begin{equation}
    \sigma_{\text{pred}} = \sqrt{\frac{1}{T-1} \sum_{i=1}^{T} \left(f\!\left(\mathbf{x};\, \bm{w},\, \mathbf{d}_i\right) - \hat{y}\right)^2}
    \label{eq:mc_std}
\end{equation}

\noindent where $\bm{w}$ represents the deterministic weights from the final training epoch, and $\mathbf{d}_i$ is the Bernoulli mask applied at the $i$-th inference pass.

\subsection{Threat and Confidence Scoring}

To enable interpretable decision-making, model output is expressed into specific \textit{Threat} and \textit{Confidence} metrics.

The \textbf{Threat score} (0--100) combines: predicted collision probability, $P_c$ trends, inverse miss distance, and time urgency. It is given by:

\begin{equation}
    \text{Threat} = 50 \cdot \left(1 + \tanh\!\left(\alpha \cdot \log_{10} P_c^{\text{pred}} + \beta\right)\right)
    \label{eq:threat}
\end{equation}

\noindent where $\alpha$ and $\beta$ are calibrated to map the typical $P_c$ range $[10^{-10},\, 10^{-2}]$ smoothly to the 0--100 scale.

\textbf{Confidence level} (0--1) combines three signals: inverse MC Dropout variance, number of observed CDMs (data uncertainty), and prediction consistency across the 50 MC iterations. Each signal is normalized to a 0--1 score and combined with weights summing to 1 to produce the final confidence. At the event level, aggregation includes: maximum threat score across all predictions (worst-case), recency-weighted average confidence, and the final observed $P_c$ from the last CDM.

\subsubsection{Operational Classification}

Events are categorized using binary thresholds of Threat~$=50$ and Confidence~$=0.5$, resulting in a four-quadrant risk classification:

\begin{center}
\begin{tabular}{llll}
\toprule
\textbf{Quadrant} & \textbf{Threat} & \textbf{Confidence} & \textbf{Interpretation} \\
\midrule
\textbf{ACT NOW}       & $\geq 50$ & $\geq 0.5$ & Immediate operational attention required \\
\textbf{WATCH CLOSELY} & $\geq 50$ & $< 0.5$    & High threat, low certainty — monitor further \\
\textbf{SAFELY IGNORE} & $< 50$    & $\geq 0.5$ & Routine conjunction — safely deprioritize \\
\textbf{NOT PRIORITY}  & $< 50$    & $< 0.5$    & Low risk with insufficient data \\
\bottomrule
\end{tabular}
\end{center}

% ============================================================
\section{Results \& Discussion}

This section presents the performance evaluation and operational analysis of the bidirectional GRU model, which was trained on 185,511 conjunction data messages collected during the first half of 2024. The results are organized into six categories: model training performance and convergence behavior, validation of the self-supervised learning approach, operational risk classification outcomes, uncertainty quantification using Monte Carlo Dropout, architectural design rationale, and identified limitations with proposed directions for future research. The test set consisted of 9,637 samples from 6,411 unique conjunction events, enabling a robust assessment of the model's generalization to previously unseen satellite encounters.

\subsection{Model Training Performance}

\begin{figure}[H]
    \centering
    \includegraphics[width=0.85\textwidth]{fig2_loss_curves.png}
    \caption{Training and Validation Loss Curves.}
    \label{fig:loss_curves}
\end{figure}

The bidirectional GRU model underwent training for 48 epochs with early stopping, selecting epoch 28 as optimal based on validation loss. Training was completed in \textbf{24.8 minutes} using standard CPU hardware. Figures~\ref{fig:loss_curves}--\ref{fig:metrics} illustrate training convergence, with smooth learning curves and no evidence of instability or divergence.

\begin{itemize}
    \item The weighted MSE loss (Figure~\ref{fig:loss_curves}) decreased from approximately 88 to 73 on the validation set.
    \item The overall MAE (Figure~\ref{fig:mae_curves}) improved consistently, indicating minimal overfitting.
    \item The collision probability prediction error (Figure~\ref{fig:pc_error}) decreased from 0.54 to 0.43 on the training set and from 0.44 to 0.41 on the validation set, demonstrating effective learning of the most operationally significant feature.
\end{itemize}

On the held-out test set (9,637 samples from 6,411 unique conjunction events):

\begin{center}
\begin{tabular}{lc}
\toprule
\textbf{Metric} & \textbf{Value} \\
\midrule
Overall MAE (all normalized features)   & 0.464 \\
Collision probability MAE               & \textbf{0.403} \\
$\log_{10}(P_c)$~$R^2$                  & \textbf{0.649} \\
Validation $P_c$~MAE                    & 0.41 \\
Test $P_c$~MAE                          & 0.404 \\
\bottomrule
\end{tabular}
\end{center}

The consistency between validation (0.41) and test (0.404) collision probability MAE confirms appropriate generalization and absence of overfitting.

\begin{figure}[H]
    \centering
    \includegraphics[width=0.85\textwidth]{fig3_mae_curves.png}
    \caption{Training and Validation MAE Curves.}
    \label{fig:mae_curves}
\end{figure}

\begin{figure}[H]
    \centering
    \includegraphics[width=0.85\textwidth]{fig4_pc_error.png}
    \caption{Collision Probability Prediction Error.}
    \label{fig:pc_error}
\end{figure}

\begin{figure}[H]
    \centering
    \includegraphics[width=0.85\textwidth]{fig5_metrics_comparison.png}
    \caption{Metrics Comparison.}
    \label{fig:metrics}
\end{figure}

\subsection{Self-Supervised Learning Validation}

A primary challenge in this study was the absence of ground-truth collision outcomes for supervised training, given the rarity of actual satellite collisions. Analyzing the relationship between model-derived threat scores and actual collision probabilities shows whether the self-supervised approach identified meaningful risk patterns or arbitrary mappings.

Figure~\ref{fig:threat_vs_pc} shows a strong positive correlation between threat scores and final $P_c$ values, with a \textbf{linear regression slope of 2.86}, confirming that the model independently discovered risk relationships from CDM sequence patterns without explicit labeling.

\begin{figure}[H]
    \centering
    \includegraphics[width=0.82\textwidth]{fig6_threat_vs_pc.png}
    \caption{Threat Score vs.\ Actual Collision Probability.}
    \label{fig:threat_vs_pc}
\end{figure}

The self-supervised framework was particularly effective at managing the extensive dynamic range of collision probabilities, spanning \textbf{eight orders of magnitude} (from approximately $10^{-10}$ to $10^{-2}$) in the dataset. By learning to predict subsequent CDM parameters from historical sequences, the model implicitly captured the evolution of credible threats over time. Events with stable, low $P_c$ trajectories received low threat scores, whereas events with increasing $P_c$ or high final values were assigned higher scores. This emergent behavior supports the core hypothesis that trajectory patterns, rather than instantaneous threshold crossings, offer superior risk discrimination.

The model's threat assessment mechanism was developed solely through minimizing prediction error on CDM sequences, without reliance on human-defined criteria for identifying ``high-risk'' trajectories. The strong correlation between threat scores and actual $P_c$ (Figure~\ref{fig:threat_vs_pc}), together with the $\log_{10}(P_c)$ $R^2$ of 0.649, demonstrates that self-supervised learning effectively extracted the underlying physics and operational logic present in the CDM data. This methodology differs fundamentally from threshold-based screening systems, which apply fixed rules irrespective of trajectory context, and it avoids the labeling bias inherent in supervised classification where human operators define threat categories.

\subsection{Operational Risk Classification}

\begin{figure}[H]
    \centering
    \includegraphics[width=0.90\textwidth]{fig7_risk_quadrant.png}
    \caption{Risk Assessment Quadrant Classification Dashboard.}
    \label{fig:risk_quadrant}
\end{figure}

The four-quadrant risk classification framework produced operationally realistic event distributions across the test set of 2,003 unique events:

\begin{center}
\begin{tabular}{llcl}
\toprule
\textbf{Quadrant} & \textbf{Events} & \textbf{Percentage} & \textbf{Description} \\
\midrule
\textbf{SAFELY IGNORE} & 1,132 & 56.5\% & Routine conjunctions — safely deprioritize \\
\textbf{NOT PRIORITY}  & 641   & 32.0\% & Low risk, insufficient observational data \\
\textbf{WATCH CLOSELY} & 149   & 7.4\%  & High-threat — further observation required \\
\textbf{ACT NOW}       & 81    & 4.0\%  & Immediate operational attention \\
\bottomrule
\end{tabular}
\end{center}

The \textbf{ACT NOW} quadrant exhibited $P_c$ values ranging from $1.46 \times 10^{-8}$ to $4.41 \times 10^{-5}$, with threat scores exceeding 50 and confidence levels above 0.5.

\begin{table}[H]
\caption{Top 20 High-Priority Events}
\label{tab:top20}
\centering
\setlength{\tabcolsep}{4pt}
\small
\begin{tabular}{lccccrcl}
\toprule
\textbf{Event ID} & \textbf{Threat} & \textbf{Conf.} & \textbf{Final $P_c$} & \textbf{TCA} & \textbf{CDMs} & \textbf{Object 1} & \textbf{Object 2} \\
\midrule
49402\_13771 & 69.9 & 0.52 & 2.08E-06 & 4/30/2024 14:08 &  5 & KOSEN-1          & SL-3 R/B           \\
32162\_43643 & 57.3 & 0.51 & 1.99E-06 & 5/27/2024 14:40 &  2 & FENGYUN 1C DEB   & YAOGAN-32 B        \\
25865\_43924 & 57.2 & 0.53 & 1.46E-08 & 2/12/2024  0:19 &  2 & SL-16 DEB        & IRIDIUM 168        \\
58598\_58597 & 55.8 & 0.52 & 2.40E-06 & 1/4/2024  10:55 &  2 & STARLINK-31079   & STARLINK-31087     \\
54073\_53841 & 55.8 & 0.82 & 2.04E-07 & 6/10/2024  4:09 & 13 & STARLINK-5146    & STARLINK-4750      \\
53918\_51891 & 55.4 & 0.85 & 1.57E-06 & 2/3/2024   5:19 & 21 & STARLINK-5012    & STARLINK-3549      \\
36861\_40399 & 55.4 & 0.55 & 3.49E-07 & 4/2/2024   8:07 &  2 & METEOR 2-8 DEB   & DMSP 5D-2 F13 DEB  \\
36878\_36270 & 55.2 & 0.51 & 7.24E-06 & 4/27/2024 21:13 &  4 & METEOR 2-7 DEB   & FENGYUN 1C DEB     \\
17642\_41508 & 54.3 & 0.50 & 4.41E-05 & 4/28/2024 22:56 &  4 & COSMOS 1275 DEB  & NOAA 16 DEB        \\
42339\_39657 & 54.3 & 0.50 & 3.83E-05 & 4/1/2024   9:18 &  3 & NOAA 16 DEB      & COSMOS 1867 COOL.  \\
52000\_53544 & 54.2 & 0.85 & 6.79E-06 & 3/6/2024   1:27 & 16 & STARLINK-3653    & STARLINK-4303      \\
45550\_48105 & 54.1 & 0.61 & 4.05E-07 & 5/22/2024  6:06 &  5 & STARLINK-1390    & STARLINK-2440      \\
52691\_53998 & 53.9 & 0.66 & 4.76E-07 & 2/3/2024  22:42 &  6 & STARLINK-3945    & STARLINK-5128      \\
30197\_21420 & 53.8 & 0.52 & 1.00E-10 & 4/13/2024 15:41 &  2 & FENGYUN 1C DEB   & SL-8 DEB           \\
50176\_53844 & 53.7 & 0.83 & 2.10E-05 & 3/4/2024  11:24 & 35 & STARLINK-3293    & STARLINK-4745      \\
54186\_52850 & 53.5 & 0.79 & 3.88E-06 & 1/15/2024  1:15 &  9 & STARLINK-5194    & STARLINK-4187      \\
53912\_51883 & 53.5 & 0.78 & 8.68E-06 & 2/16/2024 16:12 &  9 & STARLINK-5059    & STARLINK-3576      \\
46753\_45228 & 53.5 & 0.80 & 5.71E-06 & 2/11/2024 20:45 & 23 & STARLINK-1922    & STARLINK-1222      \\
47398\_45409 & 53.4 & 0.52 & 4.45E-07 & 2/11/2024 20:31 &  2 & STARLINK-2120    & STARLINK-1282      \\
46129\_46791 & 53.4 & 0.75 & 2.72E-06 & 3/4/2024   5:11 & 19 & STARLINK-1623    & STARLINK-1933      \\
\bottomrule
\end{tabular}
\end{table}

Notably, high confidence was distributed across the entire threat spectrum rather than being concentrated in low-risk events, confirming that the model distinguishes between ``confident it is safe'' and ``confident it is dangerous.'' This classification scheme enables a \textbf{96\% reduction} in urgent alerts requiring immediate analyst attention (from 2,003 events to 81 in the ACT NOW quadrant), demonstrating significant potential for operational efficiency gains in space traffic management workflows.

\subsection{Uncertainty Quantification and Confidence Calibration}

Monte Carlo Dropout, implemented through 50 stochastic forward passes during inference with dropout layers active (\texttt{training=True}), provided built-in uncertainty quantification without additional architectural complexity or ensemble training. The prediction variance across these passes naturally captures epistemic uncertainty arising from limited or ambiguous observational data.

Analysis of the test set revealed a mean uncertainty of 0.127, ranging from 0.037 to 9.81, indicating that the model appropriately expresses varying degrees of confidence across different conjunction scenarios.

\begin{figure}[H]
    \centering
    \includegraphics[width=0.82\textwidth]{fig8_confidence_vs_obs.png}
    \caption{Confidence Level vs.\ Amount of Observation Data.}
    \label{fig:confidence_obs}
\end{figure}

As illustrated in Figure~\ref{fig:confidence_obs}, confidence levels show a clear positive relationship with CDM count:
\begin{itemize}
    \item Events observed through 15--40 CDMs: mean confidence exceeding 0.7
    \item Events observed through 2--5 CDMs: confidence typically below 0.4
    \item Overall mean confidence: 0.58 across all test events
\end{itemize}

This emergent behavior — resulting solely from prediction variance patterns rather than explicit programming — demonstrates that the model naturally recognizes when it lacks sufficient information for reliable assessment. The overall mean confidence of 0.58 across all test events suggests appropriate epistemic humility, as the model avoids overconfident predictions when faced with limited or noisy data.

Importantly, the distribution of high-confidence predictions spans the full threat spectrum, from confident assessments of safe passages to confident identification of dangerous encounters. Events in the ACT NOW quadrant represent cases where both high risk and high certainty align, providing operators with clear actionable intelligence. In contrast, the WATCH CLOSELY quadrant flags potentially dangerous situations where prediction uncertainty warrants additional observations or refined orbit determination before committing to collision avoidance maneuvers. This confidence-aware classification framework addresses a fundamental limitation of traditional threshold-based screening systems, which provide binary alerts without quantifying the reliability of underlying predictions.

\subsection{Model Architecture Selection and Advantages}

The bidirectional GRU architecture was selected for its ability to capture the full temporal context in CDM sequences while maintaining computational efficiency. In contrast to unidirectional models, which process sequences in a single temporal direction, the bidirectional configuration processes each sequence in both directions. This approach enables the model to utilize comprehensive contextual information when predicting subsequent CDM parameters. Such a design is particularly valuable for conjunction risk assessment, where both early trajectory indicators and proximity to the time of closest approach influence the evolution of collision probability.

GRU cells were chosen over LSTM units due to their simplified gating mechanism, which reduces parameter count and training time, while preserving comparable sequence modeling capabilities. The resulting architecture, comprising \textbf{244,171 trainable parameters}, achieves an effective balance: it is sufficiently complex to learn nuanced trajectory patterns across the eight-order-of-magnitude $P_c$ range, yet remains compact enough to be trained in under 25 minutes on standard CPU hardware without specialized computing resources.

The self-supervised learning framework provides significant advantages compared to traditional threshold-based screening systems:

\begin{center}
\begin{tabular}{lll}
\toprule
\textbf{Aspect} & \textbf{Traditional Threshold} & \textbf{This Approach} \\
\midrule
Threat identification & Fixed $P_c$ thresholds & Learned from CDM trajectory evolution \\
Context awareness     & None — applied uniformly & Trajectory-aware discrimination \\
Labeling requirement  & Human-defined categories & No labels required \\
Inference speed       & Fast (rule lookup) & Millisecond-scale \\
Training cost         & N/A & 25 min on CPU \\
\bottomrule
\end{tabular}
\end{center}

\subsection{Limitations and Future Directions}

Several limitations of the current approach should be acknowledged:

\begin{enumerate}
    \item \textbf{Absence of ground truth collisions:} The lack of actual collision outcomes prevents direct validation, as satellite collisions are extremely rare (approximately once per year worldwide). Consequently, model performance is assessed using surrogate metrics such as collision probability predictions and trajectory evolution patterns.

    \item \textbf{Object-agnostic treatment:} The current implementation treats all space objects equally, without accounting for object-specific characteristics such as maneuverability, operational status, or spacecraft type. Active satellites capable of collision avoidance maneuvers exhibit different risk profiles compared to derelict debris.

    \item \textbf{External factors not directly modeled:} Although the model implicitly captures environmental effects through CDM sequence history, it does not directly incorporate atmospheric density variations, solar activity cycles, or gravitational perturbations.
\end{enumerate}

Future research could address these limitations through:
\begin{itemize}
    \item Incorporating \textbf{object metadata} (maneuverability indicators, cross-sectional area, operational status) as additional input features.
    \item Integrating \textbf{physical environment models} (atmospheric density forecasts, space weather indices, gravitational harmonics).
    \item Employing \textbf{ensemble methods} combining BiGRU, BiLSTM, and Transformers for improved robustness.
    \item Developing \textbf{active learning frameworks} to prioritize low-confidence events requiring refined orbit determination or supplementary observations.
\end{itemize}

% ============================================================
\section{Conclusion}

This research addresses the significant challenge of differentiating credible satellite conjunction threats from low-quality alerts amid rapidly increasing space traffic. Analysis of over 185,000 conjunction data messages collected over six months in 2024 demonstrates that self-supervised deep learning effectively learns collision risk patterns from CDM trajectory evolution, without the need for labeled collision outcomes or predefined threat categories.

The bidirectional GRU architecture achieved:
\begin{itemize}
    \item \textbf{MAE of 0.403} in collision probability prediction
    \item \textbf{$R^2 = 0.649$} (65\% of variance in log-scaled $P_c$ explained)
    \item Effective risk assessment across an eight-order-of-magnitude $P_c$ range ($10^{-10}$ to $10^{-2}$)
    \item \textbf{Slope = 2.86} in the threat score vs.\ actual $P_c$ linear regression
\end{itemize}

The operational four-quadrant classification framework yielded:
\begin{itemize}
    \item \textbf{56.5\%} safely ignorable routine conjunctions
    \item \textbf{32\%} low-priority events with insufficient data
    \item \textbf{7.4\%} high-threat cases requiring further observation
    \item \textbf{4\%} necessitating immediate operational attention
\end{itemize}

This scheme enables a \textbf{96\% reduction} in urgent alerts requiring immediate analyst review, focusing human expertise on the 81 highest-priority events from a test set of 2,003 conjunctions.

This methodology introduces several key innovations to conjunction screening:

\begin{enumerate}
    \item \textbf{Self-supervised learning} avoids labeling bias and overcomes the context insensitivity of fixed-threshold systems.
    \item \textbf{Integrated uncertainty quantification} via MC Dropout provides confidence estimates without added architectural complexity.
    \item \textbf{Trajectory-based prediction} captures temporal evolution patterns that instantaneous threshold screening cannot identify.
    \item \textbf{Computational efficiency} — 25-minute training on standard CPU hardware and millisecond-scale inference — ensures accessibility without specialized infrastructure.
\end{enumerate}

Future research should incorporate object-specific characteristics such as maneuverability and operational status, integrate physical environment models including atmospheric density and space weather, and develop active learning frameworks to optimize ground-based tracking resource allocation. Despite current limitations, this study establishes \textbf{self-supervised deep learning with uncertainty quantification} as a viable and promising approach for assessing the credibility of conjunction alerts in the evolving space traffic management landscape.

% ============================================================
\appendix
\section{Trajectory Reports for Top 20 High-Priority Events}

Detailed trajectory reports for each of the top 20 high-priority events identified in Section~3.3 (ACT NOW quadrant) are provided below.

\begin{figure}[H]
    \centering
    \includegraphics[width=\textwidth]{appendix_fig1_trajectory_reports.png}
    \caption{Trajectory reports for top-priority ACT NOW events (Part 1).}
    \label{fig:appendix1}
\end{figure}

\begin{figure}[H]
    \centering
    \includegraphics[width=\textwidth]{appendix_fig2_trajectory_reports.png}
    \caption{Trajectory reports for top-priority ACT NOW events (Part 2).}
    \label{fig:appendix2}
\end{figure}

% ============================================================
\begin{thebibliography}{10}

\bibitem{ref1}
``ESA Space Environment Report 2024,'' European Space Agency (ESA).
\newline\url{https://www.esa.int/Space_Safety/Space_Debris/ESA_Space_Environment_Report_2024}
\newline (accessed Nov.\ 10, 2025).

\bibitem{ref2}
L.~Sanchez, M.~Vasile, and E.~Minisci,
``On the Use of Machine Learning and Evidence Theory to Improve Collision Risk Management,''
\textit{2nd IAA International Conference on Space Situational Awareness},
Washington, D.C., USA, 14--16 January 2020.
(accessed Oct.\ 25, 2025).

\bibitem{ref3}
A.~K.~Mashiku, L.~K.~Newman, and D.~E.~Highsmith,
``NASA Conjunction Assessment Risk Analysis (CARA) Compendium for Artificial Intelligence and Machine Learning for Satellite Collision Avoidance,''
in \textit{Proc.\ 26th AMOS Advanced Maui Optical and Space Surveillance Technologies Conference},
Wailea, HI, USA, 2025.
[Online]. Available: \url{https://ntrs.nasa.gov/api/citations/20250008251/downloads/AMOS_2025_AIML_Paper_UpdatedContractorAddress.pdf}
(accessed Dec.\ 12, 2025).

\bibitem{ref4}
Y.~Qiao, H.-M.~Xu, W.-J.~Zhou, B.~Peng, B.~Hu, and X.~Guo,
``A BiGRU joint optimized attention network for recognition of drilling conditions,''
\textit{Petroleum Science}, vol.~20, no.~6, pp.~3624--3637, 2023.
doi: \href{https://doi.org/10.1016/j.petsci.2023.05.021}{10.1016/j.petsci.2023.05.021}
(accessed Nov.\ 3, 2025).

\bibitem{ref5}
D.~Xu \textit{et al.},
``A survey on multi-output learning,''
arXiv:1901.00248. [Online]. Available: \url{https://arxiv.org/abs/1901.00248}
(accessed Dec.\ 1, 2025).

\bibitem{ref6}
J.~Terven, D.-M.~Cordova-Esparza, J.-A.~Romero-Gonz{\'a}lez, A.~Ram{\'i}rez-Pedraza, and E.~A.~Ch{\'a}vez-Urbiola,
``A comprehensive survey of loss functions and metrics in Deep Learning,''
\textit{Artificial Intelligence Review}, Springer, 2025.
[Online]. Available: \url{https://link.springer.com/article/10.1007/s10462-025-11198-7}
(accessed Nov.\ 5, 2025).

\bibitem{ref7}
Y.~Gal and Z.~Ghahramani,
``Dropout as a Bayesian approximation: Representing model uncertainty in Deep Learning,''
arXiv:1506.02142. [Online]. Available: \url{https://arxiv.org/abs/1506.02142}
(accessed Nov.\ 29, 2025).

\bibitem{ref8}
M.~Hasan, A.~Khosravi, I.~Hossain, A.~Rahman, and S.~Nahavandi,
``Controlled dropout for uncertainty estimation,''
arXiv:2205.03109. [Online]. Available: \url{https://arxiv.org/abs/2205.03109}
(accessed Nov.\ 25, 2025).

\bibitem{ref9}
A.~Pim and T.~Pryer,
``Surrogate modelling of proton dose with Monte Carlo dropout uncertainty quantification,''
arXiv:2509.18155. [Online]. Available: \url{https://arxiv.org/abs/2509.18155}
(accessed Dec.\ 2, 2025).

\bibitem{ref10}
X.~Cao, Y.~Xu, and X.~Yang,
``Customer lifetime value prediction with uncertainty estimation using Monte Carlo Dropout,''
arXiv:2411.15944. [Online]. Available: \url{https://arxiv.org/abs/2411.15944}
(accessed Nov.\ 28, 2025).

\end{thebibliography}

\bigskip
\noindent\rule{\textwidth}{0.4pt}\\[0.3em]
\noindent\textit{Saudi Space Agency $\mid$ DebriSolver Competition $\mid$ Space Debris Conference 2026}

\end{document}
